% main.tex

% https://tex.stackexchange.com/a/748667
\unprotect
    %%% Defined similarly to \obeyMPboxdepth, \ignoreMPboxdepth,
    %%% \obeyMPboxorigin, and \normalMPboxdepth
    \permanent\protected\def\strutMPboxdepth{%
        \let\meta_relocate_box\meta_strut_box_depth%
    }

    \protected\def\meta_strut_box_depth{%
        \setbox\b_meta_graphic\vpack{%
            \box\b_meta_graphic%
            %%% Add some depth to the box without affecting its size
            \nohrule height-\openstrutdepth width 0pt depth \openstrutdepth\relax%
        }
    }

    %%% The pagecolumns output routine alredy adds a \vfil elsewhere, so the one
    %%% here isn't needed and disrupts the depth added by \strutMPboxdepth.
    \let\page_otr_fill_page_unchecked\relax
\protect

%%% Comment the line below to confirm that an overfull vbox is generated;
%%% uncomment it to confirm that the overfull vbox message goes away.
\strutMPboxdepth
\setupbodyfont[pagella]


\definepapersize
  [youtubeHD]
  [width=1920bp,height=1080bp]
\setuppapersize[youtubeHD][youtubeHD]
\setuplayout[
    location=middle,
    width=middle,
    height=middle,
    topspace=0pt,
    backspace=0pt,
    header=0pt,
    footer=0pt
]
\defineframed[myframed][
    width=800pt,
    align=normal,
    offset=1cm,
    frame=off,
    indenting={yes,medium}
]
\setupbodyfont[modern,24pt]
\ctxloadluafile{parametric.lua}
\processMPfigurefile{register_mp_cmd.mp}

\startMPinclusions
def render_theorem_box =
numeric w, h;
w := CurrentWidth;
h := CurrentHeight;

path p;
p := (-w/2, -h/2) -- (w/2, -h/2) -- (w/2, h/2) -- (-w/2, h/2) -- cycle;
setbounds currentpicture to p;

picture theorem;
theorem := textext("\getbuffer[BufA]");

pair center_anchor, topleft_anchor;
center_anchor := center theorem;
topleft_anchor := ulcorner theorem;

enddef;


def jasper_thin =
    withpen pencircle scaled .1cm
enddef ;

def jasper_thick =
    withpen pencircle scaled .2cm
enddef ;

def jasper_lightgray =
    withcolor (0.6)
enddef ;

def jasper_lightgraypenthin =
    jasper_lightgray withpen pencircle scaled .1cm
enddef ;

def jasper_lightgraypenthick =
    jasper_lightgray withpen pencircle scaled .2cm
enddef ;


def jasper_darkgray =
    withcolor (0.4)
enddef ;

def jasper_darkgraypenthin =
    jasper_darkgray withpen pencircle scaled .15cm
enddef ;

def jasper_darkgraypenthick =
    jasper_darkgray withpen pencircle scaled .3cm
enddef ;
\stopMPinclusions
\starttext
    \startbuffer[BufA]
        \startframed[myframed]
            {\bf AXIOM 2} 
            \startitemize[r][conversion=romannumerals,stopper=,left=(,right=),distance=1em]
                \item Each segment is congruent to itself: 
                \im{a=a} (reflexivity);
                \item If one segment is congruent to another, then
                the second segment is congruent to the first one:
                \im{a=b\implies b=a} (symmetry);
                \item If each of two segments is congruent to 
                a third segment, then they are also congruent 
                to each other:
                \startformula
                    \startmathalignment[n=1,align=left,location=packed,fences=sesac]
                        \NC a=c \NR
                        \NC b=c \NR
                    \stopmathalignment
                    \implies a=b
                \stopformula
                (transitivity).
            \stopitemize
            \color[white]{Segments are said to be congruent or equal
            to each other if they do not differ from one another,
            except by their location and orientation. The relation 
            of congruence is denoted by the \quotation{\im{=}} sign;
            for instance, \im{AB=MN} or \im{a=b}. The relation of 
            congruence of segments is  a fundamental undefined
            notion that obeys AXIOMS 2 and 3.}
        \stopframed
    \stopbuffer
    % Display the theorem initially, in a prominent location.
    \dorecurse{7*30}{%
        \startMPpage
            render_theorem_box ;
            numeric t;
            t := 0;
            pair anchor;
            anchor := (1 - t) * center_anchor + t * topleft_anchor + t * (-2cm, -9cm);
            label(theorem, (0,3cm) - anchor);
        \stopMPpage%
    }
    % Translate the theorem to the top right corner.
    \dorecurse{2*30+1}{%
        \startMPpage
            render_theorem_box ;
            numeric t;
            t := (#1-1)/60;
            pair anchor;
            anchor := (1 - t) * center_anchor + t * topleft_anchor + t * (-2cm, -9cm);
            label(theorem, (0,3cm) - anchor);
        \stopMPpage%
    }
    % Pause once the theorem is in place.
    \dorecurse{4*30}{%
        \startMPpage
            render_theorem_box ;
            numeric t;
            t := 1;
            pair anchor;
            anchor := (1 - t) * center_anchor + t * topleft_anchor + t * (-2cm, -9cm);
            label(theorem, (0,3cm) - anchor);
        \stopMPpage%
    }
    \startMPinclusions
        def mythm =
            render_theorem_box ;
            numeric t;
            t := 1;
            pair anchor;
            anchor := (1 - t) * center_anchor + t * topleft_anchor + t * (-2cm, -9cm);
            label(theorem, (0,3cm) - anchor);
        enddef ;
        runscript(
            "
                jasper.math.cx = -15
                jasper.math.cy = 5
            "
        ) ;
        def points_A =
            jasper_point[
                fx="cx-7"
                ,fy="cy"
            ] ;
            jasper_point[
                fx="cx"
                ,fy="cy"
            ] ;
        enddef ;
    \stopMPinclusions
    \definecolor[gray06][s=0.6]
    \startbuffer[BufA]
        \startframed[myframed]
            {\bf AXIOM 2} 
            \startitemize[r][conversion=romannumerals,stopper=,left=(,right=),distance=1em]
                \color[gray06]{\item Each segment is congruent to itself: 
                \im{a=a} (reflexivity);}
                \item If one segment is congruent to another, then
                the second segment is congruent to the first one:
                \im{a=b\implies b=a} (symmetry);
                \item If each of two segments is congruent to 
                a third segment, then they are also congruent 
                to each other:
                \startformula
                    \startmathalignment[n=1,align=left,location=packed,fences=sesac]
                        \NC a=c \NR
                        \NC b=c \NR
                    \stopmathalignment
                    \implies a=b
                \stopformula
                (transitivity).
            \stopitemize
            \color[white]{Segments are said to be congruent or equal
            to each other if they do not differ from one another,
            except by their location and orientation. The relation 
            of congruence is denoted by the \quotation{\im{=}} sign;
            for instance, \im{AB=MN} or \im{a=b}. The relation of 
            congruence of segments is  a fundamental undefined
            notion that obeys AXIOMS 2 and 3.}
        \stopframed
    \stopbuffer
    % highlight first item
    \dorecurse{4*30}{%
        \startMPpage
            mythm ;
            jasper_render_segments ;
        \stopMPpage%
    }
    % draw two points
    \dorecurse{4*30}{%
        \startMPpage
            mythm ;
            points_A ;
            jasper_render_segments ;
        \stopMPpage%
    }
    % connect the points
    \dorecurse{2*30}{%
        \startMPpage
            mythm ;
            points_A ;
            jasper_curve[
                ustart="0"
                ,ustop="1"
                ,usamples="2"
                ,fx="cx-u*7*#1/60"
                ,fy="cy"
                ,fz="0"
                ,drawcommands="jasper_thin"
            ] ;
            jasper_render_segments ;
        \stopMPpage%
    }
    % pause
    \dorecurse{2*30}{%
        \startMPpage
            mythm ;
            points_A ;
            jasper_curve[
                ustart="0"
                ,ustop="1"
                ,usamples="2"
                ,fx="cx-u*7"
                ,fy="cy"
                ,fz="0"
                ,drawcommands="jasper_thin"
            ] ;
            jasper_render_segments ;
        \stopMPpage%
    }
    % add gray line on top
    \dorecurse{2*30}{%
        \startMPpage
            mythm ;
            jasper_curve[
                ustart="0"
                ,ustop="1"
                ,usamples="2"
                ,fx="cx-u*7"
                ,fy="cy"
                ,fz="0"
                ,drawcommands="jasper_lightgraypenthin"
            ] ;
            jasper_point[
                fx="cx-7"
                ,fy="cy"
                ,drawcommands="jasper_lightgray"
            ] ;
            jasper_point[
                fx="cx"
                ,fy="cy"
                ,drawcommands="jasper_lightgray"
            ] ;
            jasper_render_segments ;
        \stopMPpage%
    }
    % rotate it around cx 360 degrees, counterclockwise
    \dorecurse{3*30}{%
        \startMPpage
            mythm ;
            points_A ;
            jasper_curve[
                ustart="0"
                ,ustop="1"
                ,usamples="2"
                ,fx="cx-u*7"
                ,fy="cy"
                ,fz="0"
                ,drawcommands="jasper_thin"
            ] ;

            jasper_curve[
                ustart="0"
                ,ustop="1"
                ,usamples="2"
                ,fx="cx+u*7*cos(pi-tau*#1/90)"
                ,fy="cy+u*7*sin(pi-tau*#1/90)"
                ,fz="1"
                ,drawcommands="jasper_lightgraypenthin"
            ] ;
            jasper_point[
                fx="cx"
                ,fy="cy"
                ,fz="1"
                ,drawcommands="jasper_lightgray"
            ] ;
            jasper_point[
                fx="cx+7*cos(pi-tau*#1/90)"
                ,fy="cy+7*sin(pi-tau*#1/90)"
                ,fz="1"
                ,drawcommands="jasper_lightgray"
            ] ;
            jasper_render_segments ;
        \stopMPpage%
    }
    % pause
    \dorecurse{2*30}{%
        \startMPpage
            mythm ;
            jasper_curve[
                ustart="0"
                ,ustop="1"
                ,usamples="2"
                ,fx="cx-u*7"
                ,fy="cy"
                ,fz="0"
                ,drawcommands="jasper_lightgraypenthin"
            ] ;
            jasper_point[
                fx="cx-7"
                ,fy="cy"
                ,drawcommands="jasper_lightgray"
            ] ;
            jasper_point[
                fx="cx"
                ,fy="cy"
                ,drawcommands="jasper_lightgray"
            ] ;
            jasper_render_segments ;
        \stopMPpage%
    }
    % pause after removing gray line
    \dorecurse{2*30}{%
        \startMPpage
            mythm ;
            points_A ;
            jasper_curve[
                ustart="0"
                ,ustop="1"
                ,usamples="2"
                ,fx="cx-u*7"
                ,fy="cy"
                ,fz="0"
                ,drawcommands="jasper_thin"
            ] ;
            jasper_render_segments ;
        \stopMPpage%
    }
    % remove segments
    \dorecurse{3*30}{%
        \startMPpage
            mythm ;
        \stopMPpage%
    }
    \startbuffer[BufA]
        \startframed[myframed]
            {\bf AXIOM 2} 
            \startitemize[r][conversion=romannumerals,stopper=,left=(,right=),distance=1em]
                \item Each segment is congruent to itself: 
                \im{a=a} (reflexivity);
                \color[gray06]{\item If one segment is congruent to another, then
                the second segment is congruent to the first one:
                \im{a=b\implies b=a} (symmetry);}
                \item If each of two segments is congruent to 
                a third segment, then they are also congruent 
                to each other:
                \startformula
                    \startmathalignment[n=1,align=left,location=packed,fences=sesac]
                        \NC a=c \NR
                        \NC b=c \NR
                    \stopmathalignment
                    \implies a=b
                \stopformula
                (transitivity).
            \stopitemize
            \color[white]{Segments are said to be congruent or equal
            to each other if they do not differ from one another,
            except by their location and orientation. The relation 
            of congruence is denoted by the \quotation{\im{=}} sign;
            for instance, \im{AB=MN} or \im{a=b}. The relation of 
            congruence of segments is  a fundamental undefined
            notion that obeys AXIOMS 2 and 3.}
        \stopframed
    \stopbuffer
    % \item If one segment is congruent to another, then
    % the second segment is congruent to the first one:
    % \im{a=b\implies b=a} (symmetry);
    \dorecurse{4*30}{%
        \startMPpage
            mythm ;
        \stopMPpage%
    }
    % outline two lines
    \dorecurse{2*30}{%
        \startMPpage
            mythm ;
            jasper_point[
                fx="cx-3"
                ,fy="cy"
            ] ;
            jasper_point[
                fx="cx-3-5"
                ,fy="cy"
            ] ;

            jasper_point[
                fx="cx+3"
                ,fy="cy"
            ] ;
            jasper_point[
                fx="cx+3+5"
                ,fy="cy"
            ] ;
            jasper_render_segments ;
        \stopMPpage%
    }
    % draw them
    \dorecurse{2*30}{%
        \startMPpage
            mythm ;
            jasper_point[
                fx="cx-3"
                ,fy="cy"
            ] ;
            jasper_point[
                fx="cx-3-5"
                ,fy="cy"
            ] ;
            jasper_curve[
                ustart="0"
                ,ustop="1"
                ,usamples="2"
                ,fx="cx-3-5*u*#1/60"
                ,fy="cy"
                ,fz="1"
                ,drawcommands="jasper_thin"
            ] ;

            jasper_point[
                fx="cx+3"
                ,fy="cy"
            ] ;
            jasper_point[
                fx="cx+3+5"
                ,fy="cy"
            ] ;
            jasper_curve[
                ustart="0"
                ,ustop="1"
                ,usamples="2"
                ,fx="cx+3+5*u*#1/60"
                ,fy="cy"
                ,fz="1"
                ,drawcommands="jasper_thin"
            ] ;
            jasper_render_segments ;
        \stopMPpage%
    }
    \startMPinclusions
        def basetwo =
            jasper_point[
                fx="cx-3"
                ,fy="cy"
            ] ;
            jasper_point[
                fx="cx-3-5"
                ,fy="cy"
            ] ;
            jasper_curve[
                ustart="0"
                ,ustop="1"
                ,usamples="2"
                ,fx="cx-3-5*u"
                ,fy="cy"
                ,fz="1"
                ,drawcommands="jasper_thin"
            ] ;

            jasper_point[
                fx="cx+3"
                ,fy="cy"
            ] ;
            jasper_point[
                fx="cx+3+5"
                ,fy="cy"
            ] ;
            jasper_curve[
                ustart="0"
                ,ustop="1"
                ,usamples="2"
                ,fx="cx+3+5*u"
                ,fy="cy"
                ,fz="1"
                ,drawcommands="jasper_thin"
            ] ;
        enddef ;
    \stopMPinclusions%
    % pause
    \dorecurse{2*30}{%
        \startMPpage
            mythm ;
            basetwo ;
            jasper_render_segments ;
        \stopMPpage%
    }
    % highlight them
    \dorecurse{2*30}{%
        \startMPpage
            mythm ;
            jasper_point[
                fx="cx-3"
                ,fy="cy"
                ,fz="5"
                ,drawcommands="jasper_lightgray"
            ] ;
            jasper_point[
                fx="cx-3-5"
                ,fy="cy"
                ,fz="5"
                ,drawcommands="jasper_lightgray"
            ] ;
            jasper_curve[
                ustart="0"
                ,ustop="1"
                ,usamples="2"
                ,fx="cx-3-5*u"
                ,fy="cy"
                ,fz="5"
                ,drawcommands="jasper_lightgraypenthin"
            ] ;

            jasper_point[
                fx="cx+3"
                ,fy="cy"
                ,fz="5"
                ,drawcommands="jasper_lightgray"
            ] ;
            jasper_point[
                fx="cx+3+5"
                ,fy="cy"
                ,fz="5"
                ,drawcommands="jasper_lightgray"
            ] ;
            jasper_curve[
                ustart="0"
                ,ustop="1"
                ,usamples="2"
                ,fx="cx+3+5*u"
                ,fy="cy"
                ,fz="5"
                ,drawcommands="jasper_lightgraypenthin"
            ] ;
            jasper_render_segments ;
        \stopMPpage%
    }
    % rotate them
    \dorecurse{3*30}{%
        \startMPpage
            mythm ;
            basetwo ;
            jasper_point[
                fx="cx-3"
                ,fy="cy"
                ,fz="5"
                ,drawcommands="jasper_lightgray"
                ,transformation="
                    matrix_multiply(
                        matrix_multiply(translate(-cx,-cy,0),zrotation(pi*#1/90))
                        ,translate(cx,cy,0)
                    )
                "
            ] ;
            jasper_point[
                fx="cx-3-5"
                ,fy="cy"
                ,fz="5"
                ,drawcommands="jasper_lightgray"
                ,transformation="
                    matrix_multiply(
                        matrix_multiply(translate(-cx,-cy,0),zrotation(pi*#1/90))
                        ,translate(cx,cy,0)
                    )
                "
            ] ;
            jasper_curve[
                ustart="0"
                ,ustop="1"
                ,usamples="2"
                ,fx="cx-3-5*u"
                ,fy="cy"
                ,fz="5"
                ,drawcommands="jasper_lightgraypenthin"
                ,transformation="
                    matrix_multiply(
                        matrix_multiply(translate(-cx,-cy,0),zrotation(pi*#1/90))
                        ,translate(cx,cy,0)
                    )
                "
            ] ;

            jasper_point[
                fx="cx+3"
                ,fy="cy"
                ,fz="5"
                ,drawcommands="jasper_lightgray"
                ,transformation="
                    matrix_multiply(
                        matrix_multiply(translate(-cx,-cy,0),zrotation(pi*#1/90))
                        ,translate(cx,cy,0)
                    )
                "
            ] ;
            jasper_point[
                fx="cx+3+5"
                ,fy="cy"
                ,fz="5"
                ,drawcommands="jasper_lightgray"
                ,transformation="
                    matrix_multiply(
                        matrix_multiply(translate(-cx,-cy,0),zrotation(pi*#1/90))
                        ,translate(cx,cy,0)
                    )
                "
            ] ;
            jasper_curve[
                ustart="0"
                ,ustop="1"
                ,usamples="2"
                ,fx="cx+3+5*u"
                ,fy="cy"
                ,fz="5"
                ,drawcommands="jasper_lightgraypenthin"
                ,transformation="
                    matrix_multiply(
                        matrix_multiply(translate(-cx,-cy,0),zrotation(pi*#1/90))
                        ,translate(cx,cy,0)
                    )
                "
            ] ;
            jasper_render_segments ;
        \stopMPpage%
    }
    % pause
    \dorecurse{3*30}{%
        \startMPpage
            mythm ;
            basetwo ;
            jasper_point[
                fx="cx-3"
                ,fy="cy"
                ,fz="5"
                ,drawcommands="jasper_lightgray"
                ,transformation="
                    matrix_multiply(
                        matrix_multiply(translate(-cx,-cy,0),zrotation(pi))
                        ,translate(cx,cy,0)
                    )
                "
            ] ;
            jasper_point[
                fx="cx-3-5"
                ,fy="cy"
                ,fz="5"
                ,drawcommands="jasper_lightgray"
                ,transformation="
                    matrix_multiply(
                        matrix_multiply(translate(-cx,-cy,0),zrotation(pi))
                        ,translate(cx,cy,0)
                    )
                "
            ] ;
            jasper_curve[
                ustart="0"
                ,ustop="1"
                ,usamples="2"
                ,fx="cx-3-5*u"
                ,fy="cy"
                ,fz="5"
                ,drawcommands="jasper_lightgraypenthin"
                ,transformation="
                    matrix_multiply(
                        matrix_multiply(translate(-cx,-cy,0),zrotation(pi))
                        ,translate(cx,cy,0)
                    )
                "
            ] ;

            jasper_point[
                fx="cx+3"
                ,fy="cy"
                ,fz="5"
                ,drawcommands="jasper_lightgray"
                ,transformation="
                    matrix_multiply(
                        matrix_multiply(translate(-cx,-cy,0),zrotation(pi))
                        ,translate(cx,cy,0)
                    )
                "
            ] ;
            jasper_point[
                fx="cx+3+5"
                ,fy="cy"
                ,fz="5"
                ,drawcommands="jasper_lightgray"
                ,transformation="
                    matrix_multiply(
                        matrix_multiply(translate(-cx,-cy,0),zrotation(pi))
                        ,translate(cx,cy,0)
                    )
                "
            ] ;
            jasper_curve[
                ustart="0"
                ,ustop="1"
                ,usamples="2"
                ,fx="cx+3+5*u"
                ,fy="cy"
                ,fz="5"
                ,drawcommands="jasper_lightgraypenthin"
                ,transformation="
                    matrix_multiply(
                        matrix_multiply(translate(-cx,-cy,0),zrotation(pi))
                        ,translate(cx,cy,0)
                    )
                "
            ] ;
            jasper_render_segments ;
        \stopMPpage%
    }
    % remove lines, pause
    \dorecurse{3*30}{%
        \startMPpage
            mythm ;
        \stopMPpage%
    }
    \startbuffer[BufA]
        \startframed[myframed]
            {\bf AXIOM 2} 
            \startitemize[r][conversion=romannumerals,stopper=,left=(,right=),distance=1em]
                \item Each segment is congruent to itself: 
                \im{a=a} (reflexivity);
                \item If one segment is congruent to another, then
                the second segment is congruent to the first one:
                \im{a=b\implies b=a} (symmetry);
                \color[gray06]{\item If each of two segments is congruent to 
                a third segment, then they are also congruent 
                to each other:
                \startformula
                    \startmathalignment[n=1,align=left,location=packed,fences=sesac]
                        \NC a=c \NR
                        \NC b=c \NR
                    \stopmathalignment
                    \implies a=b
                \stopformula
                (transitivity).}
            \stopitemize
            \color[white]{Segments are said to be congruent or equal
            to each other if they do not differ from one another,
            except by their location and orientation. The relation 
            of congruence is denoted by the \quotation{\im{=}} sign;
            for instance, \im{AB=MN} or \im{a=b}. The relation of 
            congruence of segments is  a fundamental undefined
            notion that obeys AXIOMS 2 and 3.}
        \stopframed
    \stopbuffer
    % highlight third  item of the axiom
    \dorecurse{3*30}{%
        \startMPpage
            mythm ;
        \stopMPpage%
    }
    \startMPinclusions
        numeric ir, our ;
        ir := 2 ;
        our := 5 ;
        runscript("jasper.math.ir = " & decimal(ir)) ;
        runscript("jasper.math.our = " & decimal(our)) ;
    \stopMPinclusions%
    % extend three lines
    \dorecurse{3*30}{%
        \startMPpage
            mythm ;
            numeric timer ;
            timer := #1/90 ;
            runscript("jasper.math.timer = " & decimal(timer)) ;
            for i=0 upto 2:
                runscript("jasper.math.i = " & decimal(i)) ;
                jasper_curve[
                    ustart="0"
                    ,ustop="1"
                    ,usamples="2"
                    ,fx="cx+(ir+u*our*timer)*cos(i*tau/3)"
                    ,fy="cy+(ir+u*our*timer)*sin(i*tau/3)"
                    ,fz="0"
                    ,drawcommands="jasper_thin"
                ] ;
            endfor ;
            jasper_render_segments ;
        \stopMPpage%
    }
    % pause once extended
    \dorecurse{3*30}{%
        \startMPpage
            mythm ;
            numeric timer ;
            timer := 1 ;
            runscript("jasper.math.timer = " & decimal(timer)) ;
            for i=0 upto 2:
                runscript("jasper.math.i = " & decimal(i)) ;
                jasper_curve[
                    ustart="0"
                    ,ustop="1"
                    ,usamples="2"
                    ,fx="cx+(ir+u*our*timer)*cos(i*tau/3)"
                    ,fy="cy+(ir+u*our*timer)*sin(i*tau/3)"
                    ,fz="0"
                    ,drawcommands="jasper_thin"
                ] ;
            endfor ;
            jasper_render_segments ;
        \stopMPpage%
    }
    % move the i=0 spoke into the i=1 spoke (it will eventuall move backward to the i=2 spoke)
    \dorecurse{3*30}{%
        \startMPpage
            mythm ;
            numeric timer ;
            timer := 1 ;
            runscript("jasper.math.timer = " & decimal(timer)) ;
            numeric timertwo ;
            timertwo := #1/90 ;
            runscript("jasper.math.timertwo = " & decimal(timertwo)) ;
            for i=0 upto 2:
                runscript("jasper.math.i = " & decimal(i)) ;
                if i=0:
                    jasper_curve[
                        ustart="0"
                        ,ustop="1"
                        ,usamples="2"
                        ,fx="cx+(ir+u*our*timer)*cos((i+timertwo)*tau/3)"
                        ,fy="cy+(ir+u*our*timer)*sin((i+timertwo)*tau/3)"
                        ,fz="0"
                        ,drawcommands="jasper_thin"
                    ] ;
                else:
                    jasper_curve[
                        ustart="0"
                        ,ustop="1"
                        ,usamples="2"
                        ,fx="cx+(ir+u*our*timer)*cos(i*tau/3)"
                        ,fy="cy+(ir+u*our*timer)*sin(i*tau/3)"
                        ,fz="0"
                        ,drawcommands="jasper_thin"
                    ] ;
                fi ;
            endfor ;
            jasper_render_segments ;
        \stopMPpage%
    }
    % pause before going backwards to i=2
    \dorecurse{1*30}{%
        \startMPpage
            mythm ;
            numeric timer ;
            timer := 1 ;
            runscript("jasper.math.timer = " & decimal(timer)) ;
            numeric timertwo ;
            timertwo := 1 ;
            runscript("jasper.math.timertwo = " & decimal(timertwo)) ;
            for i=0 upto 2:
                runscript("jasper.math.i = " & decimal(i)) ;
                if i=0:
                    jasper_curve[
                        ustart="0"
                        ,ustop="1"
                        ,usamples="2"
                        ,fx="cx+(ir+u*our*timer)*cos((i+timertwo)*tau/3)"
                        ,fy="cy+(ir+u*our*timer)*sin((i+timertwo)*tau/3)"
                        ,fz="0"
                        ,drawcommands="jasper_thin"
                    ] ;
                else:
                    jasper_curve[
                        ustart="0"
                        ,ustop="1"
                        ,usamples="2"
                        ,fx="cx+(ir+u*our*timer)*cos(i*tau/3)"
                        ,fy="cy+(ir+u*our*timer)*sin(i*tau/3)"
                        ,fz="0"
                        ,drawcommands="jasper_thin"
                    ] ;
                fi ;
            endfor ;
            jasper_render_segments ;
        \stopMPpage%
    }
    % move backwards to i=2
    \dorecurse{3*30}{%
        \startMPpage
            mythm ;
            numeric timer ;
            timer := 1 ;
            runscript("jasper.math.timer = " & decimal(timer)) ;
            numeric timertwo ;
            timertwo := 1 ;
            runscript("jasper.math.timertwo = " & decimal(timertwo)) ;
            numeric timerthree ;
            timerthree := #1/90 ;
            runscript("jasper.math.timerthree = " & decimal(timerthree)) ;
            for i=0 upto 2:
                runscript("jasper.math.i = " & decimal(i)) ;
                if i=0:
                    jasper_curve[
                        ustart="0"
                        ,ustop="1"
                        ,usamples="2"
                        ,fx="cx+(ir+u*our*timer)*cos((i+timertwo-2*timerthree)*tau/3)"
                        ,fy="cy+(ir+u*our*timer)*sin((i+timertwo-2*timerthree)*tau/3)"
                        ,fz="0"
                        ,drawcommands="jasper_thin"
                    ] ;
                else:
                    jasper_curve[
                        ustart="0"
                        ,ustop="1"
                        ,usamples="2"
                        ,fx="cx+(ir+u*our*timer)*cos(i*tau/3)"
                        ,fy="cy+(ir+u*our*timer)*sin(i*tau/3)"
                        ,fz="0"
                        ,drawcommands="jasper_thin"
                    ] ;
                fi ;
            endfor ;
            jasper_render_segments ;
        \stopMPpage%
    }
    % pause before returning to original position
    \dorecurse{1*30}{%
        \startMPpage
            mythm ;
            numeric timer ;
            timer := 1 ;
            runscript("jasper.math.timer = " & decimal(timer)) ;
            numeric timertwo ;
            timertwo := 1 ;
            runscript("jasper.math.timertwo = " & decimal(timertwo)) ;
            numeric timerthree ;
            timerthree := 1 ;
            runscript("jasper.math.timerthree = " & decimal(timerthree)) ;
            for i=0 upto 2:
                runscript("jasper.math.i = " & decimal(i)) ;
                if i=0:
                    jasper_curve[
                        ustart="0"
                        ,ustop="1"
                        ,usamples="2"
                        ,fx="cx+(ir+u*our*timer)*cos((i+timertwo-2*timerthree)*tau/3)"
                        ,fy="cy+(ir+u*our*timer)*sin((i+timertwo-2*timerthree)*tau/3)"
                        ,fz="0"
                        ,drawcommands="jasper_thin"
                    ] ;
                else:
                    jasper_curve[
                        ustart="0"
                        ,ustop="1"
                        ,usamples="2"
                        ,fx="cx+(ir+u*our*timer)*cos(i*tau/3)"
                        ,fy="cy+(ir+u*our*timer)*sin(i*tau/3)"
                        ,fz="0"
                        ,drawcommands="jasper_thin"
                    ] ;
                fi ;
            endfor ;
            jasper_render_segments ;
        \stopMPpage%
    }
    % return i=0 to its original position
    \dorecurse{2*30}{%
        \startMPpage
            mythm ;
            numeric timer ;
            timer := 1 ;
            runscript("jasper.math.timer = " & decimal(timer)) ;
            numeric timertwo ;
            timertwo := 1 ;
            runscript("jasper.math.timertwo = " & decimal(timertwo)) ;
            numeric timerthree ;
            timerthree := 1 ;
            runscript("jasper.math.timerthree = " & decimal(timerthree)) ;
            numeric timerfour ;
            timerfour := #1/60 ;
            runscript("jasper.math.timerfour = " & decimal(timerfour)) ;
            for i=0 upto 2:
                runscript("jasper.math.i = " & decimal(i)) ;
                if i=0:
                    jasper_curve[
                        ustart="0"
                        ,ustop="1"
                        ,usamples="2"
                        ,fx="cx+(ir+u*our*timer)*cos((i+timertwo-2*timerthree+timerfour)*tau/3)"
                        ,fy="cy+(ir+u*our*timer)*sin((i+timertwo-2*timerthree+timerfour)*tau/3)"
                        ,fz="0"
                        ,drawcommands="jasper_thin"
                    ] ;
                else:
                    jasper_curve[
                        ustart="0"
                        ,ustop="1"
                        ,usamples="2"
                        ,fx="cx+(ir+u*our*timer)*cos(i*tau/3)"
                        ,fy="cy+(ir+u*our*timer)*sin(i*tau/3)"
                        ,fz="0"
                        ,drawcommands="jasper_thin"
                    ] ;
                fi ;
            endfor ;
            jasper_render_segments ;
        \stopMPpage%
    }
    % pause after returning to the original  position
    \dorecurse{2*30}{%
        \startMPpage
            mythm ;
            numeric timer ;
            timer := 1 ;
            runscript("jasper.math.timer = " & decimal(timer)) ;
            numeric timertwo ;
            timertwo := 1 ;
            runscript("jasper.math.timertwo = " & decimal(timertwo)) ;
            numeric timerthree ;
            timerthree := 1 ;
            runscript("jasper.math.timerthree = " & decimal(timerthree)) ;
            numeric timerfour ;
            timerfour := 1 ;
            runscript("jasper.math.timerfour = " & decimal(timerfour)) ;
            for i=0 upto 2:
                runscript("jasper.math.i = " & decimal(i)) ;
                if i=0:
                    jasper_curve[
                        ustart="0"
                        ,ustop="1"
                        ,usamples="2"
                        ,fx="cx+(ir+u*our*timer)*cos((i+timertwo-2*timerthree+timerfour)*tau/3)"
                        ,fy="cy+(ir+u*our*timer)*sin((i+timertwo-2*timerthree+timerfour)*tau/3)"
                        ,fz="0"
                        ,drawcommands="jasper_thin"
                    ] ;
                else:
                    jasper_curve[
                        ustart="0"
                        ,ustop="1"
                        ,usamples="2"
                        ,fx="cx+(ir+u*our*timer)*cos(i*tau/3)"
                        ,fy="cy+(ir+u*our*timer)*sin(i*tau/3)"
                        ,fz="0"
                        ,drawcommands="jasper_thin"
                    ] ;
                fi ;
            endfor ;
            jasper_render_segments ;
        \stopMPpage%
    }
    % move i=1 and i=2 together
    \dorecurse{2*30}{%
        \startMPpage
            mythm ;
            numeric timer ;
            timer := #1/60 ;
            runscript("jasper.math.timer = " & decimal(timer)) ;
            for i=0 upto 2:
                runscript("jasper.math.i = " & decimal(i)) ;
                if i=0:
                    jasper_curve[
                        ustart="0"
                        ,ustop="1"
                        ,usamples="2"
                        ,fx="cx+(ir+u*our)*cos(i*tau/3)"
                        ,fy="cy+(ir+u*our)*sin(i*tau/3)"
                        ,fz="0"
                        ,drawcommands="jasper_thin"
                    ] ;
                elseif i=1:
                    jasper_curve[
                        ustart="0"
                        ,ustop="1"
                        ,usamples="2"
                        ,fx="cx+(ir+u*our)*cos(i*tau/3+timer*tau/6)"
                        ,fy="cy+(ir+u*our)*sin(i*tau/3+timer*tau/6)"
                        ,fz="0"
                        ,drawcommands="jasper_thin"
                    ] ;
                elseif i=2 :
                    jasper_curve[
                        ustart="0"
                        ,ustop="1"
                        ,usamples="2"
                        ,fx="cx+(ir+u*our)*cos(i*tau/3-timer*tau/6)"
                        ,fy="cy+(ir+u*our)*sin(i*tau/3-timer*tau/6)"
                        ,fz="0"
                        ,drawcommands="jasper_thin"
                    ] ;
                fi ;
            endfor ;
            jasper_render_segments ;
        \stopMPpage%
    }
    % pause after showing transitivity
    \dorecurse{5*30}{%
        \startMPpage
            mythm ;
            numeric timer ;
            timer := 1 ;
            runscript("jasper.math.timer = " & decimal(timer)) ;
            for i=0 upto 2:
                runscript("jasper.math.i = " & decimal(i)) ;
                if i=0:
                    jasper_curve[
                        ustart="0"
                        ,ustop="1"
                        ,usamples="2"
                        ,fx="cx+(ir+u*our*timer)*cos(i*tau/3)"
                        ,fy="cy+(ir+u*our*timer)*sin(i*tau/3)"
                        ,fz="0"
                        ,drawcommands="jasper_thin"
                    ] ;
                elseif i=1:
                    jasper_curve[
                        ustart="0"
                        ,ustop="1"
                        ,usamples="2"
                        ,fx="cx+(ir+u*our*timer)*cos(i*tau/3+timer*tau/6)"
                        ,fy="cy+(ir+u*our*timer)*sin(i*tau/3+timer*tau/6)"
                        ,fz="0"
                        ,drawcommands="jasper_thin"
                    ] ;
                elseif i=2 :
                    jasper_curve[
                        ustart="0"
                        ,ustop="1"
                        ,usamples="2"
                        ,fx="cx+(ir+u*our*timer)*cos(i*tau/3-timer*tau/6)"
                        ,fy="cy+(ir+u*our*timer)*sin(i*tau/3-timer*tau/6)"
                        ,fz="0"
                        ,drawcommands="jasper_thin"
                    ] ;
                fi ;
            endfor ;
            jasper_render_segments ;
        \stopMPpage%
    }
    % remove segments
    \dorecurse{2*30}{%
        \startMPpage
            mythm ;
        \stopMPpage%
    }
    \startbuffer[BufA]
        \startframed[myframed]
            {\bf AXIOM 2} 
            \startitemize[r][conversion=romannumerals,stopper=,left=(,right=),distance=1em]
                \item Each segment is congruent to itself: 
                \im{a=a} (reflexivity);
                \item If one segment is congruent to another, then
                the second segment is congruent to the first one:
                \im{a=b\implies b=a} (symmetry);
                \item If each of two segments is congruent to 
                a third segment, then they are also congruent 
                to each other:
                \startformula
                    \startmathalignment[n=1,align=left,location=packed,fences=sesac]
                        \NC a=c \NR
                        \NC b=c \NR
                    \stopmathalignment
                    \implies a=b
                \stopformula
                (transitivity).
            \stopitemize
            Segments are said to be congruent or equal
            to each other if they do not differ from one another,
            except by their location and orientation. \color[white]{The relation 
            of congruence is denoted by the \quotation{\im{=}} sign;
            for instance, \im{AB=MN} or \im{a=b}. The relation of 
            congruence of segments is  a fundamental undefined
            notion that obeys AXIOMS 2 and 3.}
        \stopframed
    \stopbuffer
    % Segments are said to be congruent or equal
    % to each other if they do not differ from one another,
    % except by their location and orientation. 
    \dorecurse{7*30}{%
        \startMPpage
            mythm ;
        \stopMPpage%
    }
    \startbuffer[BufA]
        \startframed[myframed]
            {\bf AXIOM 2} 
            \startitemize[r][conversion=romannumerals,stopper=,left=(,right=),distance=1em]
                \item Each segment is congruent to itself: 
                \im{a=a} (reflexivity);
                \item If one segment is congruent to another, then
                the second segment is congruent to the first one:
                \im{a=b\implies b=a} (symmetry);
                \item If each of two segments is congruent to 
                a third segment, then they are also congruent 
                to each other:
                \startformula
                    \startmathalignment[n=1,align=left,location=packed,fences=sesac]
                        \NC a=c \NR
                        \NC b=c \NR
                    \stopmathalignment
                    \implies a=b
                \stopformula
                (transitivity).
            \stopitemize
            Segments are said to be congruent or equal
            to each other if they do not differ from one another,
            except by their location and orientation. The relation 
            of congruence is denoted by the \quotation{\im{=}} sign;
            \color[white]{for instance, \im{AB=MN} or \im{a=b}. The relation of 
            congruence of segments is  a fundamental undefined
            notion that obeys AXIOMS 2 and 3.}
        \stopframed
    \stopbuffer
    % The relation 
    % of congruence is denoted by the \quotation{\im{=}} sign;
    \dorecurse{4*30}{%
        \startMPpage
            mythm ;
        \stopMPpage%
    }
    \startbuffer[BufA]
        \startframed[myframed]
            {\bf AXIOM 2} 
            \startitemize[r][conversion=romannumerals,stopper=,left=(,right=),distance=1em]
                \item Each segment is congruent to itself: 
                \im{a=a} (reflexivity);
                \item If one segment is congruent to another, then
                the second segment is congruent to the first one:
                \im{a=b\implies b=a} (symmetry);
                \item If each of two segments is congruent to 
                a third segment, then they are also congruent 
                to each other:
                \startformula
                    \startmathalignment[n=1,align=left,location=packed,fences=sesac]
                        \NC a=c \NR
                        \NC b=c \NR
                    \stopmathalignment
                    \implies a=b
                \stopformula
                (transitivity).
            \stopitemize
            Segments are said to be congruent or equal
            to each other if they do not differ from one another,
            except by their location and orientation. The relation 
            of congruence is denoted by the \quotation{\im{=}} sign;
            for instance, \im{AB=MN} or \im{a=b}. \color[white]{The relation of 
            congruence of segments is  a fundamental undefined
            notion that obeys AXIOMS 2 and 3.}
        \stopframed
    \stopbuffer
    % for instance, \im{AB=MN} or \im{a=b}.
    \dorecurse{3*30}{%
        \startMPpage
            mythm ;
        \stopMPpage%
    }
    \startbuffer[BufA]
        \startframed[myframed]
            {\bf AXIOM 2} 
            \startitemize[r][conversion=romannumerals,stopper=,left=(,right=),distance=1em]
                \item Each segment is congruent to itself: 
                \im{a=a} (reflexivity);
                \item If one segment is congruent to another, then
                the second segment is congruent to the first one:
                \im{a=b\implies b=a} (symmetry);
                \item If each of two segments is congruent to 
                a third segment, then they are also congruent 
                to each other:
                \startformula
                    \startmathalignment[n=1,align=left,location=packed,fences=sesac]
                        \NC a=c \NR
                        \NC b=c \NR
                    \stopmathalignment
                    \implies a=b
                \stopformula
                (transitivity).
            \stopitemize
            Segments are said to be congruent or equal
            to each other if they do not differ from one another,
            except by their location and orientation. The relation 
            of congruence is denoted by the \quotation{\im{=}} sign;
            for instance, \im{AB=MN} or \im{a=b}. The relation of 
            congruence of segments is  a fundamental undefined
            notion that obeys AXIOMS 2 and 3.
        \stopframed
    \stopbuffer
    % The relation of 
    % congruence of segments is  a fundamental undefined
    % notion that obeys AXIOMS 2 and 3.
    \dorecurse{15*30}{%
        \startMPpage
            mythm ;
        \stopMPpage%
    }
\stoptext 
